\documentclass[10pt]{article}

\usepackage[utf8]{inputenc}

\usepackage[T1]{fontenc}

\usepackage{amsmath}

\usepackage{amsfonts}

\usepackage{amssymb}

\usepackage[version=4]{mhchem}

\usepackage{stmaryrd}

\usepackage{bbold}

\usepackage{CJKutf8}



\begin{document}

группы $G$. Для локально компактных групп П. к. групшы $G$ принято называть коммутативное кольцо, определенное операциями прямой суммы и тензорного произведения в множестве классов эквивалентности непрөрывных унитарных представлений группы $G$. Изучение структуры П. к. плодотворно для компактных групп, где оно приводит к теории двойственности в терминах блок-алгебр, а также в более общем случае для групп типа I, где изучение структуры П. к. может быть сведено к изучению структуры тензорных произведений неприводимых унитарных представлений.



А. И. Штерн.



ПРЕДСТАВЛЕНИИ ТЕОРИЯ - теория, изучающая гомоморфизмы полугрупп (в частности, групп), алгебр или других алгебраич. систем в соответствующие системы эндоморфизмов нек-рой подходящей структуры. Особенно часто рассматриваются л и в е й н ы е п р е д с т а в л е н и я, т. е. гомоморфизмы полугрупп, групा, ассоциативных алгебр или алгебр Ли в полугруппу, группу, алгебру линейных преобразований нек-рого векторного пространства $V$. Такое представление наз. также ли н е й ны м п р е дставлением в пространстве $V$, а $V$ наз. п р о с т а вст в о п ред с т а ления. Часто под П. т. понимают именно теорию линейных представлений. Если пространство $V$ конечномерно, то его размерность наз. р а 3 м е р н о с т ь ю п р е дс т а в л е н и я, а само представление - ко н е чн о м е р в ы м. Таким образом, различаются конеqномерные и бесконечномерныө представления. Представление наз. т о q н ы м, если оно инъективно.



Изучение линөйных представлений полугрупп, групп и алгебр Ли сводится к изучению линейных представлений ассоциативных алгеб $p$. А именно, линейные представления полугрупп (линейные представления групn) в пространстве $V$ над полем $k$ находятся в естественном взаимно однозначном соответствии с представлениями соответствующей полугрупповой (групповой) алгебры над $k$ в простравстве $V$. Представления алгебры Ли $L$ над $k$ взаимно однозначно соответствуют линейным представлениям өе универсальной обертывающөй алгебры.



Задание линейвого представления $\varphi$ ассоциативной алгебры $A$ в пространстве $V$ равносильно заданию на $V$ структуры $A$-модуля, называемого м о д у ле м п р е дс т а в л өн и я $\varphi$. При рассмотрении представлений группы $G$ пли алгебры Ли $L$ говорят также о $G$-модулях или $L$-модулях (см. Модуль). Гомоморфизмы модулей представлений наз. сплетаюцияи операторами. Изоморфным модулям соответствуют э к в и в а ле н тны е п р өд с т а в е ния. Подмодуль модуля $V$ представления $\varphi$ - это подпространство $W \subset V$, инвариантное относительно $\varphi$; индуцируемое в $W$ представление наз. п о д п р е д с т а в л е и и м, а представление, индуцируемое в фактормодуле $\mathrm{V} / \mathrm{W}$, фа к т о п п е д с т а в л е н и е п представления $\varphi$. Прямые суммы модулөй соответствуют прямым суммам представлений, неразложимые модули - неразложимым представлениям, простые модули - неприводимым представлениям, а полупростые модули - вполне приводимым представлениям. Определяются также төнзорное произведение линейных предетавлений, внешняя и симметрич. степени представления (см. Тензорное произведение представлений).



Наряду с абстрактной (или алгебраической) П. т. существует теория представлений топологич. объектов, напр. топологич. групп или банаховых алгебр (см. Непрерывное представление, Представление топологической гpyппы).



ПРЕДСТАВЛЕНИЯ КЛАССИЧЕСКИХ ГРУІІ в т е в з о р а х - линейные представления групп

$G L(V), \quad S L(V), \quad O(V, f), \quad S O(V, f), \quad S p(V, f) \quad$ (где $V$ есть $n$-мерное векторное пространство над полөм $k$, $f$ - невырожденная симметрическая или знакопере. менная билинейная форма на $V$ ) в инвариантных подпространствах тензорных степеней $T^{m}(V)$ пространства $V$. Если $k$ - поле нулевой характеристики, то все веприводимые полиномиальные линейные представления указанных групп реализуются в тензорах.



В случае $k=\mathbb{C}$ переqисленные выше группы являются комплексными группами Ли. Для всех них, кроме $G L(V)$, все (дифференцируемые) линейные представления полиномиальны; всякое линейное представление группы $G L(V)$ имеет вид $g \mapsto(\operatorname{det} g)^{k} R(g)$, где $k \in \mathbb{Z}$, а $R$ - полиномиальное линейное представление. Классические компактные группы Ли $U_{n}, S U_{n}$, $O_{n}, S O_{n}, S p_{n}$ имеют те же комплексные линейные представления и те же инвариантные подпространства в пространствах тензоров, что и их комплексные оболочки $G L_{n}(\mathbb{C}), \quad S L_{n}(\mathbb{C}), \quad O_{n}(\mathbb{C}), \quad S O_{n}(\mathbb{C}), \quad S p_{2 n}(\mathbb{C})$. Поэтому результаты теории линейных прөдставлений, полученные для классических комплексных групп Ли, переносятся на соответствующие компактные группы и наоборот («унитарный трюо) Вейля). В частности, с помощью интегрирования по компактной группе доказывается полная приводимость линейвых представлений классических комплексных групп Ли.



Естественное линейное представление группы $G L(V)$ в пространстве $T^{m}(V)$ определяется по формуле $g\left(v_{1} \otimes \ldots \otimes v_{m}\right)=g v_{i} \bigotimes \ldots \otimes g v_{m}, g \in G L(V), v_{i} \in V$.



В том же пространстве определено линейное представление симметрич. группы $S_{m}$ :



$\sigma\left(v_{1} \otimes \ldots \otimes v_{m}\right)=v_{\sigma^{-1}(1)} \otimes \ldots \otimes v_{\sigma^{-1}(m)}, \sigma \in S_{m}, v_{i} \in V$. Операторы этих двух представлений перестановочны; тем самым в $T^{m}(V)$ определено линейное представление группы $G L(V) \times S_{m}$. Если char $k=0$, то простравство $T^{m}(V)$ может быть разложено в прямую сумму минимальных $\left(G L(V) \times S_{m}\right)$-инвариантных подпространств:



\[

T^{m}(V)=\Sigma_{\lambda} V_{\lambda} \otimes U_{\lambda} .

\]



Здесь суммирование происходит по всем разбиениям $\lambda$ числа $m$, содержащим не более $n$ слагаемых, $U_{\lambda}-$ простравство абсолютно неприводимого представления $T_{\lambda}$ группы $S_{m}$, отвечающего разбиению $\lambda$ (см. Представление симметрической группы), а V\begin{CJK}{UTF8}{mj}入\end{CJK}-пространство нек-рого абсслютно неприводимого представления $R_{\lambda}$ группы $G L(V)$. Разбиения $\lambda$ удобно предстөвлять в виде наборов $\left(\lambda_{1}, \lambda_{2}, \ldots, \lambda_{n}\right)$ целых неотрицательных чисел, удовлетворяющих условиям $\lambda_{1} \geqslant \lambda_{2} \geqslant \ldots \geqslant \lambda_{n}$, $\Sigma_{i} \lambda_{i}=m$.



Подпространство $V_{\lambda} \otimes U_{\lambda} \subset T^{m}(V)$ разлагается в сумму минимальных $G L(V)$-инвариантных подпространств, в каждом из к-рых реализуется представление $R_{\lambda}$. Эти подпространства могут быть явно получены применением к $T^{m}(V)$ Юнга симлетризаторов, связавных с разбиением $\lambda$. Напр., для разбиения $\lambda=(m, 0, \ldots, 0)$ (соответственно $\lambda=(1, \ldots, 1,0, \ldots, 0)$ при $m \leqslant n)$ $\operatorname{dim} U_{\lambda}=1$ и $V_{\lambda} \otimes U_{\lambda}$ есть минимальное $G L(V)$-инвариантное подпространство, состоящее из всех симметрич. (соответственно кососимметрич.) тензоров.



Представление $R_{\lambda}$ характеризуется следующими свойствами. Пусть $B \subset G L(V)$ - подгруппа, состоящая из линейных операторов, к-рыө в нек-ром базисе $\left(e_{1}, \ldots\right.$, $\left.e_{n}\right)$ пространства $V$ записываются верхними треугольными матрицами. Тогда операторы $R_{\lambda}(b), b \in B$, имеют единственный (с точностью до числового множителя) общий собственный вектор $v_{\lambda}$, называемый с т а р шा и м в е к т о р м представления $R_{\lambda}$. Соответствующее собственное значение (с т а $\mathrm{p}$ ш й в е с представления $R_{\lambda}$ ) равно $b_{11}^{\lambda_{1}} \ldots b_{n n}^{\lambda_{n}}$, где $b_{i i}$ есть $i$-й диагональный әлемент матрицы оператора $b$ в базисе $\left(e_{1}, \ldots, e_{n}\right)$.





\end{document}